%% ch18 --- Why you need chaos: a field guide to prediction
%%
%% Stub, generated by tools/gen_stubs.py from tools/chapters.json.
%% The brief below is the contract for this chapter: what it must argue, what
%% it must buy the reader, and what it must NOT take from its neighbours.
%% Writing the chapter means deleting the stub block below (the one that opens
%% on the next \chapter line) and replacing it with the chapter text -- after
%% which this file is never regenerated.
%% If the brief turns out to be wrong once the code has been run, change it in
%% the manifest and record the contradiction in CLAUDE.md under *Resolved
%% questions*.

\chapter{Why you need chaos: a field guide to prediction}
\label{ch:why}

\chapterstub{%
Argument: the book's conclusion, stated as a practical skill rather than a
sentiment. Knowing chaos theory changes what you do: you ask for the horizon
of any prediction, you prefer an ensemble to a point, you separate questions
about paths from questions about statistics, you stop attributing big
effects to small causes after the fact, you distrust patterns in short
irregular data, and you check that a simulation's answer survives a change
of step size and precision. The chapter collects the book's measurements
into one table of horizons (logistic map, Lorenz, Henon, double pendulum,
the solar system as quoted) and turns them into seven questions to ask of
any forecast, model or claim \dash{} each question citing the chapter where
the reader measured it. It also says plainly what chaos theory is NOT: not a
theory of everything, not complexity or disorder in general, not an excuse
for fatalism (the climate and the control chapters say why), and not a
licence to call anything unpredictable butterfly-shaped. The reader leaves
able to explain to a colleague, with numbers they computed, why some things
can be predicted for a century and others not for a fortnight. Payoff: the
field guide. Traps to elicit (the last catalogue check): unpredictable means
unknowable; chaos means anything goes; small causes always matter. Listings:
the horizon table recomputed from the book's own values. Labs: apply the
seven questions to a model the reader writes. Figures: the map of
predictability \dash{} horizon against error for every system in the book on
one plot.

\medskip
\textbf{Word budget:} 3000.
}
